\section{The contraction margin at the first horizons}\label{sec:contraction}

The decomposition makes $V_{\omega,L}$ computable from elementary functions.
This section builds it, says what its contraction data must be to second order
in $\omega$, and reports the computation at $L=\log3,\log5,\log7$ and beyond.
The assembly is elementary enough to run in \emph{double precision}, so the
core of it is a registered check programme replayed with the manuscript; the
long horizons, where the margin is $10^{-17}$ and smaller, stay exploratory at
40--90 digits. Both are described in Appendix~\ref{app:numerics}, and every
number below has a record beside it.

\subsection{The assembly}\label{sec:assembly}

In the orthonormal sine basis $e_j(x)=\sqrt{2/L}\sin(j\pi(x+\frac L2)/L)$,
$j=1,\dots,N$, on $I_L$, the Galerkin matrix of the compression is
\begin{equation}\label{eq:galerkin}
 V_{jk}=\ip{e_j}{V_{\omega,L}e_k}=\int_0^Lk_\omega(t)\,S_{jk}(t)\dd t,\qquad
 S_{jk}(t)=\int_{-L/2+t}^{L/2}e_j(x)\,e_k(x-t)\dd x ,
\end{equation}
because $V_{\omega,L}$ is a causal convolution. The one-sided autocorrelations
are elementary: with $\alpha=\pi/L$ and $\sigma=(-1)^{j+k}$,
\begin{equation}\label{eq:S}
\begin{aligned}
 S_{jj}(t)&=\frac{(L-t)\cos(j\alpha t)+\sin(j\alpha t)/(j\alpha)}L,\\
 S_{jk}(t)&=\frac1L\Big[\frac{\sigma\sin(k\alpha t)-\sin(j\alpha t)}{(j-k)\alpha}
 +\frac{\sigma\sin(k\alpha t)+\sin(j\alpha t)}{(j+k)\alpha}\Big]\qquad(j\neq k).
\end{aligned}
\end{equation}
The kernel is assembled in the form of Proposition~\ref{prop:factor}:
$k_\omega=\sum_{n<e^L}\ct_n\kappa_\omega(\cdot-\log n)$ with
$\kappa_\omega=k^\Gamma_\omega-4\omega b\,J_{-b}-4\omega a\,J_a$ and
$J_c(\tau)=\int_0^\tau e^{c(\tau-u)}k^\Gamma_\omega(u)\dd u$. The $t$-integral is
the only quadrature, and the whole difficulty is the spike
$k^\Gamma_\omega(\tau)\sim\omega\pi^\omega2^\omega\tau^{\omega-1}$, which carries
half its mass below $\tau=2^{-1/\omega}$ --- below $2^{-500}$ at
$\omega=0.002$. The way through is not to evaluate it but to put it in the
quadrature \emph{weight}.

\begin{proposition}[One atom is a weight times an analytic function]\label{prop:analytic}
Let $0<\omega<\frac12$, $a=\frac12-\omega$, $b=\frac12+\omega$, and
\begin{equation}\label{eq:Ghat}
 G(\tau)=\frac{2\pi^\omega}{\Gamma(\omega)}\Big(\frac{2\sinh\tau}\tau\Big)^{\omega-1}e^{\tau/2},
 \qquad
 \Psi_c(\tau)=\int_0^1v^{\omega-1}e^{c\tau(1-v)}G(\tau v)\dd v .
\end{equation}
Then $k^\Gamma_\omega(\tau)=\tau^{\omega-1}G(\tau)$, $J_c(\tau)=\tau^\omega\Psi_c(\tau)$
for every $c\in\R$, and
\begin{equation}\label{eq:atom}
 \kappa_\omega(\tau)=\tau^{\omega-1}\Gh(\tau),\qquad
 \Gh=G-4\omega\tau\big(b\,\Psi_{-b}+a\,\Psi_a\big),
\end{equation}
where $G$, $\Psi_c$ and $\Gh$ are holomorphic on the disc $|\tau|<\pi$ and
$\Gh(0)=G(0)=\pi^\omega2^\omega/\Gamma(\omega)$.
\end{proposition}
\begin{proof}
$2\sinh\tau=e^\tau(1-e^{-2\tau})$ puts $k^\Gamma_\omega$ in the stated form.
$\sinh\tau/\tau$ is entire, equals $1$ at the origin and vanishes exactly at
$\tau=i\pi k$, $k\neq0$, so $\log(2\sinh\tau/\tau)$ and each of its powers are
holomorphic on $|\tau|<\pi$. Substituting $u=\tau v$ in $J_c$ and extracting
$\tau^{\omega-1}$ from $k^\Gamma_\omega(\tau v)$ gives the $\Psi$ form;
$\Psi_c$ is holomorphic there because $|\tau v|\leq|\tau|$ on $v\in[0,1]$ and
the $v$-integral converges absolutely.
\end{proof}

Every integral of the assembly is therefore
$\int_0^T\tau^{\omega-1}(\text{holomorphic})\dd\tau$ with $T<\pi$ at every
horizon used here, and the rule for it is \emph{Gauss--Jacobi for the weight
$\tau^{\omega-1}$}, whose nodes and weights come from the closed-form monic
Jacobi recurrence with $\alpha=0$, $\beta=\omega-1$ by Golub--Welsch. Three
things follow. The local coordinate is automatic, because each atom is
integrated over its whole support $[\log n,L)$ in its own $\tau=t-\log n$: one
rule per atom, no splitting at intermediate atoms, no global time formed by
subtraction. Nothing underflows: the smallest node is $\sim10^{-5}$, because
$\tau^{\omega-1}$ is integrated exactly by the rule rather than evaluated ---
which is what makes double precision possible. And convergence is geometric:
thirty nodes per atom give every digit that sixty do, where the substitution
$\tau=u^{1/\omega}$ followed by a tanh--sinh rule needed $262$ evaluations and
still left a floor of $6\times10^{-12}$ in $\lambda_{\min}(D)$
\cite{ContractionNote}.

The identity that certifies the assembly is elementary too.

\begin{proposition}[One atom, in closed form]\label{prop:laplace}
For $\re p>b$,
\begin{equation}\label{eq:laplace-atom}
 \int_0^\infty e^{-p\tau}k^\Gamma_\omega(\tau)\dd\tau
 =\pi^\omega\frac{\Gamma\big(\frac{a+p}2\big)}{\Gamma\big(\frac{b+p}2\big)},
 \qquad
 \int_0^\infty e^{-p\tau}\kappa_\omega(\tau)\dd\tau
 =\pi^\omega\frac{\Gamma\big(\frac{a+p}2\big)}{\Gamma\big(\frac{b+p}2\big)}
 \cdot\frac{(p+a)(p-b)}{(p+b)(p-a)} .
\end{equation}
\end{proposition}
\begin{proof}
The first integrand is
$\frac{2\pi^\omega}{\Gamma(\omega)}e^{-(p-\omega+1/2)\tau}(1-e^{-2\tau})^{\omega-1}$,
and $x=e^{-2\tau}$ turns the integral into
$\frac{\pi^\omega}{\Gamma(\omega)}\int_0^1x^{\frac{a+p}2-1}(1-x)^{\omega-1}\dd x
=\frac{\pi^\omega}{\Gamma(\omega)}B\big(\frac{a+p}2,\omega\big)$, with
$\frac{a+p}2+\omega=\frac{b+p}2$. The two smoothings multiply the transform by
$1-\frac{4\omega b}{p+b}-\frac{4\omega a}{p-a}$, whose numerator over the common
denominator is $p^2+2\omega p-ab-4\omega(a+b)p=p^2-2\omega p-ab=(p+a)(p-b)$,
since $a+b=1$ and $b-a=2\omega$.
\end{proof}

\begin{remark}[The Beta variable, explicitly]\label{rem:betavar}
Dividing \eqref{eq:laplace-atom} by its value at $p=0$ gives
$B(\frac{a+p}2,\omega)/B(\frac a2,\omega)=\E\,U^{p/2}$ with
$U\sim\mathrm{Beta}(\frac a2,\omega)$, which is \eqref{eq:betamoments}. The
change of variable in the proof says what $U$ \emph{is}: $U=e^{-2\tau}$, so the
archimedean delay is $\tau=-\frac12\log U$ and the Bessel additivity of
Proposition~\ref{prop:bessel} is a statement about $e^{-2\tau}$. It is the same
variable that Section~\ref{sec:horocycle} identifies as the horocycle
coordinate, $U=(1+|t|^2)^{-1}$, and that Section~\ref{sec:nobeta} shows no
radial chain produces.
\end{remark}

The right-hand sides of \eqref{eq:laplace-atom} contain $\Gamma$ and nothing
else: no $\zeta$, no $\xi$, no zeros. Asserting them against the assembled
kernel tests the Beta kernel, the local coordinates and the $R_\omega$
correction together, and the registered programme makes that assertion at
$p\in\{2,3,5,8\}$ and $\omega\in\{0.01,0.1\}$, agreeing to $1.4\times10^{-13}$.

Three quantities are then read from the $N\times N$ matrix: $\norm{V_N}$, which
must be below one; $\lambda_{\min}(D_N)$ with $D_N=I-V_N^{\mathsf T}V_N$,
compared through \eqref{eq:intro-firstorder} with $m_L^{(N)}$, the smallest
eigenvalue of $Q_L$ in the same basis, computed in closed form from
\eqref{eq:target} by the assembler of \cite{FrameBound}; and the first-order law
on the fixed vector $e_1$, where $Q_L[e_1]$ is thousands of times $m_L$ and the
check is insensitive to the noise floor.

\subsection{What the numbers should be}\label{sec:theory}

Write $G=G_{0,L}$ for the compressed generator at $\omega=0$ and split it into
Hermitian and skew parts, $G=Q+A$, so that $Q=Q_L$ is the Weil form and $A$ is
a Hilbert transform on the interval together with the odd part of the comb and
smooth terms.

\begin{proposition}[The defect to second order]\label{prop:secondorder}
On the compression,
\begin{equation}\label{eq:secondorder}
 D_{\omega,L}=2\omega\,Q-\omega^2\big(2Q^2+[Q,A]\big)+O(\omega^3),
\end{equation}
and if $f$ is a normalized eigenvector of $Q$ with eigenvalue $m$ then
$\ip f{D_{\omega,L}f}/2\omega=m-\omega m^2+O(\omega^2)$. In particular
$\lambda_{\min}(D_{\omega,L})/2\omega=m_L-\omega m_L^2+O(\omega^2)$: the relative
correction to the first-order law at the minimizer is $-\omega m_L$ and
nothing else.
\end{proposition}
\begin{proof}
By \eqref{eq:generator}, $K_\omega=\exp(-\int_0^\omega a_{\omega'}\dd\omega')$,
and $a_\omega$ is even in $\omega$, so $K_\omega=\exp(-\omega a_0+O(\omega^3))$
as symbols on $\re p>b$, hence as causal kernels. By
Lemma~\ref{lem:multiplicative} compression is a homomorphism, so
$V_{\omega,L}=\exp(-\omega G+O(\omega^3))=I-\omega G+\frac{\omega^2}2G^2+O(\omega^3)$
and
$V^*V=I-\omega(G+G^*)+\omega^2\big(G^*G+\frac12G^2+\frac12G^{*2}\big)+O(\omega^3)$.
With $G=Q+A$: $G^*G=Q^2+QA-AQ-A^2$ and $\frac12(G^2+G^{*2})=Q^2+A^2$, so the
bracket is $2Q^2+[Q,A]$. For the eigenvector,
$\ip f{[Q,A]f}=\ip{Qf}{Af}-\ip{A^*f}{Qf}=2\re\ip{Qf}{Af}=2m\,\re\ip f{Af}=0$ since $A^*=-A$,
and $\ip f{Q^2f}=m^2$. The eigenvalue statement is first-order perturbation
theory for the lowest eigenvalue of $Q$, which is an eigenvalue with a smooth
eigenvector (the form is of logarithmic order on the interval, bounded below
with compact resolvent) and is taken simple, as the computations show. The
expansions are of matrix elements $\int_0^Lk_\omega(t)S(t)\dd t$ against smooth
$S$, which are analytic in $\omega$: the spike integrates against a constant to
$\epsilon^\omega$.
\end{proof}

\begin{proposition}[The Galerkin defect]\label{prop:galerkin}
Let $E_N$ be the span of the first $N$ basis functions, $P_N$ its projection,
$V_N=P_NV_{\omega,L}P_N$, and $f_N$ the normalized minimizer of $Q$ on $E_N$
with eigenvalue $m_L^{(N)}$. Then
\begin{equation}\label{eq:galerkin-defect}
 D_N:=I_N-V_N^{\mathsf T}V_N=P_ND_{\omega,L}P_N+P_NV^*(I-P_N)VP_N\ \geq\ P_ND_{\omega,L}P_N ,
\end{equation}
and
\begin{equation}\label{eq:galerkin-expansion}
 \frac{\ip{f_N}{D_Nf_N}}{2\omega}=m_L^{(N)}-\omega\big(m_L^{(N)}\big)^2
 +\frac\omega2\Big(\norm{(I-P_N)Af_N}^2-\norm{(I-P_N)Qf_N}^2\Big)+O(\omega^2).
\end{equation}
If the defect is instead computed \emph{nested} in $E_{N'}\supset E_N$,
$D_{N\subset N'}:=I_N-(V_{N'}^{\mathsf T}V_{N'})_{N\times N}=P_N(I-V^*P_{N'}V)P_N$,
the term $\frac\omega2\norm{(I-P_N)Gf_N}^2$ hidden in
\eqref{eq:galerkin-expansion} is replaced by $\frac\omega2\norm{(I-P_{N'})Gf_N}^2$,
and the first-order coefficient tends, as $N'\to\infty$, to
\begin{equation}\label{eq:nested-limit}
 -\big(m_L^{(N)}\big)^2-\norm{(I-P_N)Qf_N}^2-\ip{(I-P_N)Qf_N}{(I-P_N)Af_N},
\end{equation}
which is negative as soon as the first two terms dominate.
\end{proposition}
\begin{proof}
$V_N^{\mathsf T}V_N=P_NV^*P_NVP_N$ and $P_N\leq I$ give the identity and the
inequality. On $f_N\in E_N$, $(I-P_N)Vf_N=-\omega(I-P_N)Gf_N+O(\omega^2)$, so the
extra term is $\omega^2\norm{(I-P_N)Gf_N}^2+O(\omega^3)$. In
Proposition~\ref{prop:secondorder} evaluated on $f_N$, which is an eigenvector
of $P_NQP_N$ but not of $Q$, write $Qf_N=m_L^{(N)}f_N+(I-P_N)Qf_N$; then
$\norm{Qf_N}^2=(m_L^{(N)})^2+\norm{(I-P_N)Qf_N}^2$ and
$\re\ip{Qf_N}{Af_N}=\re\ip{(I-P_N)Qf_N}{(I-P_N)Af_N}$, since $\re\ip{f_N}{Af_N}=0$. Collecting these with
$\norm{(I-P_N)(Q+A)f_N}^2$ expanded gives \eqref{eq:galerkin-expansion} and,
with $\norm{(I-P_{N'})Gf_N}\to0$, the limit \eqref{eq:nested-limit}.
\end{proof}

So a finite-basis computation must show a \emph{positive} first-order excess
over $m_L^{(N)}$, dominated by the energy of $Af_N$ outside the basis, and that
excess must not merely shrink under nesting but overshoot: the limit
\eqref{eq:nested-limit} is negative. Both are tested below.

\begin{proposition}[The first-order tail]\label{prop:tail}
For $L>0$, integrating over $[0,L)$,
\begin{equation}\label{eq:tail}
 \lim_{\omega\downarrow0}\frac1\omega\int_L^\infty k_\omega(t)\dd t
 =8\sinh\frac L2-2\sum_{n<e^L}\frac{\Lambda(n)}{\sqrt n}-\log2\pi-\gamma_E
 -\int_0^L\Big(2\gammak(u)-\frac1u\Big)\dd u-\log L .
\end{equation}
At $L=\log3$ the five terms are $4.618802$, $-0.980258$, $-1.837877$,
$-0.577216$, $-0.592935$, total $0.6305160$; at $\log5$ and $\log7$ the totals
are $0.6447529$ and $0.8355651$.
\end{proposition}
\begin{proof}
$\int_0^\infty k_\omega=K_\omega(0)=\xi(\frac12-\omega)/\xi(\frac12+\omega)=1$
for every $\omega$, so the tail is minus the head, and
$\partial_\omega|_0\int_0^Lk_\omega=\int_0^Lk_1$ with $k_1=-g_0$ of
\eqref{eq:g0}. The pole terms integrate to
$-2\int_0^L(e^{-t/2}+e^{t/2})\dd t=-8\sinh\frac L2$, the constant to $\log\pi$,
the comb to $2\sum_{n<e^L}\Lambda(n)n^{-1/2}$, and the finite part, by
\eqref{eq:fp} with $\phi=\mathbf 1_{[0,L)}$ and $E_1(x)=-\gamma_E-\log x+O(x)$,
to $\gamma_E+\log2+\int_0^L(2\gammak-\frac1u)\dd u+\log L$. Collect and negate.
\end{proof}

This checks the mass of the assembled kernel by a formula that integrates the
poles, the digamma and the primes directly, where the kernel carries them
through the Beta convolution and two exponential smoothings; it uses no zeros.

\subsection{Results}\label{sec:results}

All runs: $L=\log n$ parsed exactly with atoms taken at $\log m<L$ strictly,
Gauss--Jacobi with 40 nodes per atom, 40--90 digits. In the same basis at
$L=\log3$, $m_L^{(16)}=6.84655\times10^{-8}$, $m_L^{(24)}=6.2475140\times10^{-8}$,
$m_L^{(32)}=6.21505\times10^{-8}$ (even block); the converged value is
$5.66\times10^{-8}$ \cite{FrameBound} (recorded $5.54\times10^{-8}$
\cite{CriticalPath}), so the sine basis is far from converged at these $N$, and
every comparison is \emph{within a basis}, which is what
Proposition~\ref{prop:galerkin} is about.

\begin{table}[ht]\centering\small
\caption{The $\omega$-sweep at $L=\log3$, $N=24$. The fifth column is the
relative excess over $m_L^{(24)}$ divided by $\omega$, which
Proposition~\ref{prop:galerkin} says is constant to leading order; the last is
$(1-\norm{P_NVe_1}^2)/(2\omega\,Q_L[e_1])$ with
$Q_L[e_1]=4.3474053\times10^{-4}$.}\label{tab:sweep}
\begin{tabular}{@{}lccccc@{}}
\toprule
$\omega$ & $1-\norm{V_N}$ & $\lambda_{\min}(D_N)/2\omega$ & ratio to $m_L^{(24)}$ & (ratio$-1$)$/\omega$ & $e_1$ ratio\\
\midrule
0.002 & $1.25085\times10^{-10}$ & $6.2542610\times10^{-8}$ & 1.0010800 & 0.53998 & 0.998834\\
0.01 & $6.28139\times10^{-10}$ & $6.2813856\times10^{-8}$ & 1.0054216 & 0.54216 & 0.994676\\
0.02 & $1.26313\times10^{-9}$ & $6.3156283\times10^{-8}$ & 1.0109026 & 0.54513 & 0.990607\\
0.05 & $3.21053\times10^{-9}$ & $6.4210662\times10^{-8}$ & 1.0277794 & 0.55559 & 0.985816\\
0.1 & $6.60857\times10^{-9}$ & $6.6085743\times10^{-8}$ & 1.0577926 & 0.57793 & 1.002027\\
\bottomrule
\end{tabular}
\end{table}

\paragraph{Contraction and the first-order law.} $V_N$ is a strict contraction
at every $\omega$ (Table~\ref{tab:sweep}), with $1-\norm{V_N}=\frac12\lambda_{\min}(D_N)$
to all printed digits as it must be for $\lambda_{\min}\ll1$, and
$\lambda_{\min}(D_N)/2\omega\to m_L^{(24)}$ as $\omega\downarrow0$. The transfer
side, assembled from $k^\Gamma_\omega$, $\ct_n$ and $R_\omega$, and the form
side, assembled from \eqref{eq:target}, agree at first order in $\omega$. The
fifth column is the content of Proposition~\ref{prop:galerkin}: the excess is
linear in $\omega$ with intercept $0.5394$ and a visible $O(\omega)$ curvature,
so it is the Galerkin truncation term and not the $-\omega m_L$ of
Proposition~\ref{prop:secondorder}, which is $10^{-9}$ here. Reading the
intercept and computing $\norm{(I-P_{24})Qf_{24}}^2=0.24775\,m_L^{(24)}$ directly
from the closed-form Weil matrix in $240$ modes gives
$\norm{(I-P_{24})Af_{24}}^2=1.3265\,m_L^{(24)}$ --- the sine-tail energy of the
skew partner on the near-null vector, without ever assembling $A$.

\begin{table}[ht]\centering\small
\caption{Basis size and the nested defect at $\omega=0.1$, $L=\log3$. The
nested ratio crosses below one, as \eqref{eq:nested-limit} requires; the
computable part of that limit predicts $0.97523$.}\label{tab:nested}
\begin{tabular}{@{}lccc@{}}
\toprule
$N$ & $N'$ (nested) & ratio to $m_L^{(N)}$ & excess\\
\midrule
16 & -- & 1.06508 & $+6.5\%$\\
24 & -- & 1.05779 & $+5.8\%$\\
32 & -- & 1.04536 & $+4.5\%$\\
24 & 40 & 1.03784 & $+3.8\%$\\
24 & 64 & 1.02075 & $+2.1\%$\\
24 & 96 & 1.01028 & $+1.0\%$\\
24 & 144 & 1.00253 & $+0.3\%$\\
24 & 192 & 0.99848 & $-0.2\%$\\
\bottomrule
\end{tabular}
\end{table}

\paragraph{The nested defect.} Table~\ref{tab:nested} carries the nesting to
$N'=192$. The excess is positive, decreases with $N$ and with $N'$, and
\emph{crosses below one} between $N'=144$ and $N'=192$: the first-order
coefficient of the nested defect is heading for \eqref{eq:nested-limit}, whose
first two terms alone give a limiting ratio
$1-\omega\big(m_L^{(24)}+\norm{(I-P_{24})Qf_{24}}^2/m_L^{(24)}\big)=0.97523$ at
$\omega=0.1$. An earlier reading of two nestings as a $1/N'$ law
\cite{ContractionNote} described the first step only; measured against the
limit the approach is slower, and the apparent steepening of the raw ratios is
the crossing. At $\omega=0.05$, $N'=64$ gives $1.00637$.

\begin{table}[ht]\centering\small
\caption{The three horizons at $N=24$. Atoms are the integers below $e^L$;
$m_L^{(24)}$ is the even block. The excess over $m_L^{(24)}$ is linear in
$\omega$ at each horizon, with slope $0.54$, $1.20$, $1.85$.}\label{tab:horizons}
\begin{tabular}{@{}lccccccc@{}}
\toprule
$L$ & atoms & $m_L^{(24)}$ & $\omega$ & $1-\norm{V_N}$ & $\lambda_{\min}(D_N)/2\omega$ & ratio & $e_1$ ratio\\
\midrule
$\log3$ & 2 & $6.2475\times10^{-8}$ & 0.002 & $1.2509\times10^{-10}$ & $6.25426\times10^{-8}$ & 1.00108 & 0.998834\\
 & & & 0.01 & $6.2814\times10^{-10}$ & $6.28139\times10^{-8}$ & 1.00542 & 0.994676\\
$\log5$ & 4 & $2.5445\times10^{-17}$ & 0.002 & $5.1012\times10^{-20}$ & $2.55060\times10^{-17}$ & 1.00241 & 0.997521\\
 & & & 0.01 & $2.5750\times10^{-19}$ & $2.57497\times10^{-17}$ & 1.01199 & 0.988141\\
$\log7$ & 6 & $2.1067\times10^{-22}$ & 0.002 & $4.2290\times10^{-25}$ & $2.11449\times10^{-22}$ & 1.00372 & 0.997491\\
 & & & 0.01 & $2.1456\times10^{-24}$ & $2.14560\times10^{-22}$ & 1.01849 & 0.988228\\
\bottomrule
\end{tabular}
\end{table}

\paragraph{The horizons.} Table~\ref{tab:horizons}. The within-basis caveat
matters more here than at $\log3$: the recorded converged margins at
$\log5,\log7$ are $9.3\times10^{-18}$ and $6.8\times10^{-28}$
\cite{CriticalPath,FrameBound}, against $m_L^{(24)}=2.5\times10^{-17}$ and
$2.1\times10^{-22}$ in the $24$-mode sine basis --- a factor $3$ at $\log5$ and
six orders of magnitude at $\log7$. Both sides of every comparison below are
computed in that same truncated basis, and none of them is a statement about the
converged margin. That said, $V_{\omega,L}$ is a strict contraction at every
horizon and shift tested, by margins down to $4\times10^{-25}$; the excess over $m_L^{(24)}$ is exactly linear in $\omega$ at
each, with a slope constant to three digits across a factor of five in
$\omega$; and the margins themselves run from $6.2\times10^{-8}$ to
$2.1\times10^{-22}$. \emph{Over fourteen orders of magnitude of margin, a
kernel built from $\Gamma$, $\sinh$, $\exp$ and the integers below $e^L$
reproduces the smallest eigenvalue of the localized Weil form up to a term that
is linear in $\omega$ and computable.} Nesting at $\log5$ in $64$ modes gives
$0.98793$, again straddling.

Two horizons further out, at $\omega=0.01$ and $N=24$, the plain Galerkin
defect overstates $m_L^{(24)}$ by a factor $5.26$ at $L=\log{11}$
($m_L^{(24)}=9.358\times10^{-28}$) and $16.6$ at $L=\log{13}$
($4.849\times10^{-29}$), while nesting in $N'=64$ returns $0.98595$ and
$0.98780$. The reason is structural: the margin collapses superexponentially in
$L$, the first-order truncation term --- a tail energy in a fixed number of
sine modes --- does not, and their ratio runs $0.005,0.012,0.018$ at
$\log3,\log5,\log7$ and then $4.3$ and $15.6$. \textbf{Past $\log7$ the nested
defect is not a refinement but a requirement.}

\begin{table}[ht]\centering\small
\caption{The mass on the window at $L=\log3$, and the first-order limit at the
three horizons against Proposition~\ref{prop:tail}.}\label{tab:mass}
\begin{tabular}{@{}lcc@{}}
\toprule
$\omega$ & $\int_0^Lk_\omega$ & $(1-\int_0^Lk_\omega)/\omega$\\
\midrule
0.002 & 0.99874046831 & 0.6297658\\
0.01 & 0.99373231869 & 0.6267681\\
0.02 & 0.98753944334 & 0.6230278\\
0.05 & 0.96940739694 & 0.6118521\\
0.1 & 0.94066265818 & 0.5933734\\
\midrule
\multicolumn{3}{@{}l@{}}{\emph{extrapolated to $\omega=0$, against \eqref{eq:tail}:}}\\
$\log3$ & 0.6305153 & (closed form 0.6305160)\\
$\log5$ & 0.6447519 & (closed form 0.6447529)\\
$\log7$ & 0.8355627 & (closed form 0.8355651)\\
\bottomrule
\end{tabular}
\end{table}

\paragraph{The mass.} Table~\ref{tab:mass} against Proposition~\ref{prop:tail}:
the assembled kernel has the right first-order mass beyond the horizon, to
relative $1.2\times10^{-6}$, $1.5\times10^{-6}$ and $2.8\times10^{-6}$ at the
three horizons, by a formula that knows nothing of the assembly. The
extrapolations are two-point, in $\omega=0.002$ and $0.01$, and are limited by
their own second-order error and not by the quadrature: the masses themselves
agree between the double-precision and the 40-digit runs to $10^{-13}$.

\subsection{Reading}

The elementary assembly is validated end to end, at three horizons, and now in
a registered check. The kernel side and the form side agree at the first-order
law; their disagreement is the one Proposition~\ref{prop:galerkin} prescribes,
is linear in $\omega$ with a measured coefficient at each horizon, and
overshoots under nesting exactly as \eqref{eq:nested-limit} requires. The mass
agrees with an independent closed form at all three horizons, and the kernel
agrees with the closed Laplace transform \eqref{eq:laplace-atom} to
$10^{-13}$. There is a working instrument that computes $\norm{V_{\omega,L}}$,
$\lambda_{\min}(D_{\omega,L})$ and their $\omega$-dependence from the integers
below $e^L$ and elementary functions, at any horizon whose atoms one is willing
to sum.

What it says is what Proposition~\ref{prop:dichotomy} says it must:
contraction, with margin $\omega m_L$ to first order and a positive
second-order correction. Nothing about the zeros is learned at horizons where
$m_L>0$ is already known; the dichotomy is about \emph{all} horizons.

Proposition~\ref{prop:secondorder} is worth keeping for its own sake: the
defect of the compressed all-pass is, to second order, $2\omega Q-\omega^2(2Q^2+[Q,A])$,
so the only first-order correction to the first-order law at the minimizer is
$-\omega m_L^2$, and the commutator of the Weil form with its skew partner ---
the compressed Hilbert transform of the interval --- governs the second-order
behaviour away from the minimizer.
